\begin{figure}[htbp]
\centering
% Connectivity only. These planar drawing coordinates are not X_0.
% All circles and lettering are native TikZ and LaTeX.
\begin{tikzpicture}[
  x=1cm,y=1cm,
  nucleus/.style={circle,draw=black!65,fill=white,minimum size=8mm,
    inner sep=1pt,font=\small},
  bond/.style={draw=black!65,line width=.7pt}]
  \node[nucleus,fill=figInk,text=white] (C) at (-1.1,0) {$\mathrm C_6$};
  \node[nucleus,fill=figNitrogen,text=white] (N) at (1.1,0) {$\mathrm N_7$};
  \node[nucleus] (H1) at (-2.7,0) {$\mathrm H_1$};
  \node[nucleus] (H2) at (-1.85,1.0) {$\mathrm H_2$};
  \node[nucleus] (H3) at (-1.85,-1.0) {$\mathrm H_3$};
  \node[nucleus] (H4) at (1.85,1.0) {$\mathrm H_4$};
  \node[nucleus] (H5) at (1.85,-1.0) {$\mathrm H_5$};
  \draw[bond] (C)--(N);
  \draw[bond] (C)--(H1);
  \draw[bond] (C)--(H2);
  \draw[bond] (C)--(H3);
  \draw[bond] (N)--(H4);
  \draw[bond] (N)--(H5);
  \node[font=\small] at (-1.85,-1.65) {methyl hydrogens};
  \node[font=\small] at (1.85,-1.65) {amino hydrogens};
\end{tikzpicture}
\caption{Nuclear labels and bonds for methylamine. All five hydrogens
are like nuclei in $S$. The drawing specifies connectivity, with the
three-dimensional geometry suppressed.}
\label{fig:methylamine-labels}
\end{figure}
